\subsection{Preliminary Technical Survey}
\label{sec:survey}

A preliminary literature and technical survey has been initiated to determine whether the topic has sufficient technical depth and reliable primary material for the final term paper.

One important source identified is Samuel Naffziger, Kevin Lepak, Milam Paraschou, and Mahesh Subramony's paper, ``AMD Chiplet Architecture for High-Performance Server and Desktop Products,'' presented at ISSCC 2020 \cite{naffziger2020chiplet}. This source is particularly relevant because it discusses the architectural and engineering motivation behind AMD's chiplet approach, including the increasing difficulty and cost associated with large leading-edge dies and the limited benefit of advanced process scaling for many I/O structures.

A second source identified is the AMD Zen 2 presentation by Suggs et al. at Hot Chips 31 \cite{suggs2019zen2}. Zen 2 is important to the proposed study because the compute-die/I/O-die organization became central to both mainstream Ryzen and second-generation EPYC processors. The material provides an initial basis for studying the Core Complex Die (CCD), I/O Die (IOD), and AMD Infinity Fabric.

For the cache component, a key primary source identified is John Wuu et al.'s ISSCC 2022 paper, ``3D V-Cache: The Implementation of a Hybrid-Bonded 64MB Stacked Cache for a 7nm x86-64 CPU'' \cite{wuu2022vcache}. This work is directly relevant because it links the familiar computer-architecture problem of increasing effective cache capacity with a modern physical-integration method in which additional cache SRAM is vertically stacked over the processor rather than being implemented entirely in the planar CPU die.

AMD's technical documentation for the Zen 3 EPYC generation has also been shortlisted \cite{amd2022epyc7003}. It gives a useful architectural description of the relationship between Zen 3 cores, the Core Complex, the Core Complex Die, and the shared L3 cache.

At this stage, these references have been examined at a \textbf{preliminary level}: bibliographic details, abstracts, architectural summaries, and selected technical material. Five sources are cited in this report, but none is recorded as having received a complete close reading. A structured source-by-source reading log will be maintained during the next phase.
